\documentclass[a4paper,10pt,twoside]{article}
\usepackage{fontspec}
\setmainfont{Liberation Sans}
\setsansfont{Liberation Sans}
\renewcommand{\familydefault}{\sfdefault}
\usepackage{amsmath,amssymb}
\usepackage{geometry}
\usepackage{multicol}
\usepackage{graphicx}
\usepackage{fancyhdr}
\usepackage{lastpage}
\usepackage{booktabs}
\usepackage{listings}
\usepackage{eso-pic}
\usepackage{tikz}
\usepackage{array}
\usepackage{tabularx}
\usepackage{enumitem}
\usepackage{microtype}

\geometry{
    a4paper,
    top=1.8cm,
    bottom=1.8cm,
    left=1.2cm,
    right=1.2cm,
    headsep=14pt,
    footskip=20pt
}

\setlength{\columnsep}{16pt}
\setlength{\columnseprule}{0.35pt}

% Persistent Subtle Watermark on Every Page
\newsavebox{\watermarkbox}
\sbox{\watermarkbox}{%
    \begin{tabular}{c}
        \includegraphics[width=6.0cm]{/home/johan/GATE/GATE_CS_PYQ_DB/assets/studiva_butterfly_icon_transparent.jpeg} \\[-0.2em]
        {\fontsize{14}{17}\selectfont\bfseries Studiva}
    \end{tabular}%
}
\AddToShipoutPictureBG{%
    \begin{tikzpicture}[remember picture, overlay]
        \node[opacity=0.035] at (current page.center) {\usebox{\watermarkbox}};
    \end{tikzpicture}%
}

% Running Headers & Footers
\pagestyle{fancy}
\fancyhf{}
\fancyhead[C]{\small\textbf{Studiva GATE Question Bank}}
\fancyhead[LE]{\small\textit{Computer Science \& IT (2011--2025)}}
\fancyhead[RO]{\small\textit{Studiva Education}}
\fancyfoot[LE,RO]{\scriptsize\textbf{Studiva} $\cdot$ \texttt{www.studiva.co.in}}
\fancyfoot[C]{\small\textbf{Page \thepage\ of \pageref{LastPage}}}
\fancyfoot[RE,LO]{\scriptsize\texttt{contact@studiva.co.in}}
\renewcommand{\headrulewidth}{0.4pt}
\renewcommand{\footrulewidth}{0.4pt}

% Listings Style (Clean, professional)
\lstset{
    basicstyle=\scriptsize\ttfamily,
    numbers=none,
    frame=single,
    rulecolor=\color{black!30},
    tabsize=2,
    breaklines=true,
    showstringspaces=false
}

\begin{document}

% ==================== PAGE 1: FORMAL TITLE & COVER PAGE ====================
\thispagestyle{empty}
\vspace*{1.5cm}
\begin{center}
    \includegraphics[width=2.5cm]{/home/johan/GATE/GATE_CS_PYQ_DB/assets/studiva_butterfly_icon_transparent.jpeg}\\[1.5em]
    {\fontsize{24}{28}\selectfont\textbf{GATE COMPUTER SCIENCE \& IT}}\\[0.5em]
    {\fontsize{16}{20}\selectfont\textbf{PREVIOUS YEARS QUESTION BANK (2011 -- 2025)}}\\[0.6em]
    {\normalsize Detailed Examination Papers with Complete Official Answer Keys}\\[2.0em]
    \rule{0.65\textwidth}{0.8pt}\\[2.0em]
    {\large\textbf{Studiva GATE Examination Series}}\\[0.4em]
    {\small Comprehensive Authentic Question Bank for GATE Computer Science \& Information Technology}\\[2.8cm]

    \begin{tabular}{|ll|}
        \hline
        \textbf{Subject:} & Computer Science and Information Technology (CS/IT) \\
        \textbf{Period Covered:} & 2011 to 2025 \\
        \textbf{Examination Papers:} & 1 Complete Verified Paper Sets \\
        \textbf{Contents:} & Examination Questions \& Official Final Answer Keys \\
        \textbf{Published by:} & Studiva (\texttt{www.studiva.co.in}) \\
        \hline
    \end{tabular}

    \vfill
    {\small\textbf{Studiva $\cdot$ \texttt{www.studiva.co.in} $\cdot$ \texttt{contact@studiva.co.in}}}
\end{center}

\newpage

% ==================== PAGE 2: PUBLICATION IMPRINT & NOTES ====================
\thispagestyle{empty}
\vspace*{0.6cm}
\begin{center}
    {\large\textbf{PUBLICATION INFORMATION \& COPYRIGHT NOTICE}}\\[0.4em]
    \rule{0.6\textwidth}{0.5pt}
\end{center}

\vspace{1.2em}
\noindent\begin{tabular}{ll}
\textbf{Title:} & Studiva GATE Question Bank: Computer Science \& IT (2011--2025) \\
\textbf{Publisher:} & Studiva \\
\textbf{Website:} & \texttt{www.studiva.co.in} \\
\textbf{Official Inquiries:} & \texttt{contact@studiva.co.in} \\
\textbf{Edition:} & First Edition (2026) \\
\textbf{Syllabus Scope:} & Official GATE Syllabus (General Aptitude and CS / IT) \\
\end{tabular}

\vspace{1.8em}
\noindent\fbox{\parbox{\dimexpr\linewidth-2\fboxsep-2\fboxrule}{%
\textbf{COPYRIGHT NOTICE}\\[0.4em]
Copyright \copyright\ 2026 Studiva. All Rights Reserved.\\[0.4em]
No part of this publication may be reproduced, stored in a retrieval system, or transmitted in any form or by any means, electronic, mechanical, photocopying, recording, or otherwise, without prior written permission of Studiva, except for individual academic study.
}}

\vspace{1.5em}
\noindent\fbox{\parbox{\dimexpr\linewidth-2\fboxsep-2\fboxrule}{%
\textbf{PUBLISHER'S NOTE: ARCHIVAL QUESTION FORMAT}\\[0.4em]
This compendium has been meticulously curated and verified against official GATE question papers and final answer keys.
\par\medskip\noindent
\textbf{Notice on Image-Based Questions:}\\
Questions containing diagrams, circuit schematics, graphs, or image-based figures from the original examination papers have been omitted in this text edition due to non-readable format in historical source archives. All text-based conceptual, analytical, mathematical, algorithmic, and numerical questions are fully included with verified official keys.
}}

\vspace{1.5em}
\begin{center}
\textbf{EXAMINATION STRUCTURE \& PRACTICE GUIDELINES}
\end{center}
\begin{enumerate}[leftmargin=*]
    \item \textbf{Sectional Mark Division:} Each paper comprises General Aptitude (GA) and Computer Science (CS). Questions are clearly grouped into 1-mark and 2-mark blocks at the head of each section.
    \item \textbf{Question Types:} Includes Multiple Choice Questions (MCQ), Multiple Select Questions (MSQ), and Numerical Answer Type (NAT).
    \item \textbf{Official Answer Keys:} For self-evaluation without distraction, complete official answer keys are printed at the conclusion of each respective examination paper set.
\end{enumerate}

\newpage

% ==================== PAGE 3: TABLE OF CONTENTS ====================
\begin{center}
    {\Large\textbf{TABLE OF CONTENTS}}\\[0.4em]
    \rule{0.5\textwidth}{0.6pt}
\end{center}
\vspace{1em}

\noindent\begin{tabularx}{\linewidth}{|X|l|c|r|}
\hline
\textbf{Examination Paper} & \textbf{Session / Shift} & \textbf{Total Questions} & \textbf{Content Type} \\
\hline

GATE 2018 Computer Science & Shift CS & 58 & Questions + Answer Key \\ \hline
\end{tabularx}

\vspace{1.5em}
\noindent{\footnotesize\textit{Note: All questions in this edition have been typeset from verified digital records and cross-referenced with official answer keys released by the organising IIT/IISc institutes.}}

\newpage

% ==================== PAPER: 2018_CS ====================
\begin{center}
{\Large\textbf{GATE 2018 EXAMINATION}}\\[3pt]
{\normalsize\textbf{COMPUTER SCIENCE AND INFORMATION TECHNOLOGY}}\\[2pt]
{\small\textbf{Session / Shift: CS}}\\[5pt]
\rule{\linewidth}{0.8pt}\vspace{-4pt}
{\footnotesize\textbf{Duration: 180 Minutes} \hfill \textbf{Total Marks: 100} \hfill \textbf{Questions: 58}}\\
\vspace{-4pt}\rule{\linewidth}{0.4pt}
\end{center}
\vspace{4pt}
\noindent\fbox{\parbox{\dimexpr\linewidth-2\fboxsep-2\fboxrule}{%
\footnotesize\textbf{Examination Instructions:}\\[1pt]
1. Question Types: Multiple Choice Questions (MCQ), Multiple Select Questions (MSQ), and Numerical Answer Type (NAT).\\
2. Negative Marking Scheme: For 1-mark MCQ, $1/3$ mark is deducted; for 2-mark MCQ, $2/3$ mark is deducted. There is no negative marking for MSQ and NAT questions.\\
3. Mark Division: Questions are grouped by mark allocation within each section as specified below. Full official answer key is provided at the end of this paper.
}}
\vspace{10pt}
\begin{multicols}{2}

\vspace{8pt}
\noindent\textbf{\large\underline{GENERAL APTITUDE}}\\[2pt]
\noindent{\footnotesize\textit{[Section Block: 1-mark questions]}}
\vspace{4pt}

\noindent\textbf{Q.1} \hfill {\footnotesize\textsf{[MCQ]}}
\par\vspace{2pt}
{\bfseries ``From where are they bringing their books? \_\_\_\_\_\_\_\_ bringing \_\_\_\_\_\_\_ books from
\_\_\_\_\_.''
\par\smallskip
The words that best fill the blanks in the above sentence are \par}
\par\vspace{2pt}\noindent
\begin{tabular*}{\columnwidth}{@{\extracolsep{\fill}}p{0.46\columnwidth}p{0.46\columnwidth}@{}}
\textbf{(A)}~Their, they're, there & \textbf{(B)}~They're, their, there \\
\addlinespace[2pt]
\textbf{(C)}~There, their, they're & \textbf{(D)}~They're, there, there
\end{tabular*}
\par\vspace{6pt}\hrule height 0.25pt\vspace{6pt}
\noindent\textbf{Q.1} \hfill {\footnotesize\textsf{[MCQ]}}
\par\vspace{2pt}
{\bfseries Which one of the following is a closed form expression for the generating function of the
sequence \{an\}, where an = 2n + 3 for all n = 0, 1, 2,… ? \par}
\par\vspace{2pt}\noindent
\textbf{(A)}~2 \newline 3 \newline (1 \newline )x \newline −\par\vspace{2pt}\noindent
\textbf{(B)}~2 \newline 3 \newline (1 \newline ) \newline x \newline x \newline −\par\vspace{2pt}\noindent
\textbf{(C)}~2 \newline 2 \newline (1 \newline ) \newline x \newline x \newline − \newline −\par\vspace{2pt}\noindent
\textbf{(D)}~2 \newline 3 \newline (1 \newline ) \newline x \newline x \newline − \newline −
\par\vspace{6pt}\hrule height 0.25pt\vspace{6pt}
\noindent\textbf{Q.2} \hfill {\footnotesize\textsf{[MCQ]}}
\par\vspace{2pt}
{\bfseries ``A \_\_\_\_\_\_\_\_\_ investigation can sometimes yield new facts, but typically organized ones
are more successful.''
\par\smallskip
The word that best fills the blank in the above sentence is \par}
\par\vspace{2pt}\noindent
\begin{tabular*}{\columnwidth}{@{\extracolsep{\fill}}p{0.46\columnwidth}p{0.46\columnwidth}@{}}
\textbf{(A)}~meandering & \textbf{(B)}~timely \\
\addlinespace[2pt]
\textbf{(C)}~consistent & \textbf{(D)}~systematic
\end{tabular*}
\par\vspace{6pt}\hrule height 0.25pt\vspace{6pt}
\noindent\textbf{Q.2} \hfill {\footnotesize\textsf{[MCQ]}}
\par\vspace{2pt}
{\bfseries Consider the following C program.
\par\smallskip
\#include<stdio.h>
struct Ournode\{
char x,y,z;
\};
\par\smallskip
int main()\{
struct Ournode p = \{'1', '0', 'a'+2\};
struct Ournode *q = \&p;
printf ("\%c, \%c", *((char*)q+1), *((char*)q+2));
return 0;
\}
\par\smallskip
The output of this program is: \par}
\par\vspace{2pt}\noindent
\begin{tabular*}{\columnwidth}{@{\extracolsep{\fill}}p{0.46\columnwidth}p{0.46\columnwidth}@{}}
\textbf{(A)}~0, c & \textbf{(B)}~0, a+2 \\
\addlinespace[2pt]
\textbf{(C)}~'0', 'a+2' & \textbf{(D)}~'0', 'c'
\end{tabular*}
\par\vspace{6pt}\hrule height 0.25pt\vspace{6pt}
\noindent\textbf{Q.3} \hfill {\footnotesize\textsf{[MCQ]}}
\par\vspace{2pt}
{\bfseries The area of a square is 𝑑. What is the area of the circle which has the diagonal of the
square as its diameter? \par}
\par\vspace{2pt}\noindent
\textbf{(A)}~𝜋𝑑\par\vspace{2pt}\noindent
\textbf{(B)}~𝜋𝑑2\par\vspace{2pt}\noindent
\textbf{(C)}~1 \newline 4 𝜋𝑑2\par\vspace{2pt}\noindent
\textbf{(D)}~1 \newline 2 𝜋𝑑
\par\vspace{6pt}\hrule height 0.25pt\vspace{6pt}
\noindent\textbf{Q.4} \hfill {\footnotesize\textsf{[MCQ]}}
\par\vspace{2pt}
{\bfseries What would be the smallest natural number which when divided either by 20 or by 42 or
by 76 leaves a remainder of 7 in each case? \par}
\par\vspace{2pt}\noindent
\begin{tabular*}{\columnwidth}{@{\extracolsep{\fill}}p{0.46\columnwidth}p{0.46\columnwidth}@{}}
\textbf{(A)}~3047 & \textbf{(B)}~6047 \\
\addlinespace[2pt]
\textbf{(C)}~7987 & \textbf{(D)}~63847
\end{tabular*}
\par\vspace{6pt}\hrule height 0.25pt\vspace{6pt}
\noindent\textbf{Q.4} \hfill {\footnotesize\textsf{[MCQ]}}
\par\vspace{2pt}
{\bfseries Let $\oplus$ and ⊙ denote the Exclusive OR and Exclusive NOR operations, respectively.
Which one of the following is NOT CORRECT? \par}
\par\vspace{2pt}\noindent
\textbf{(A)}~𝑃$\oplus$𝑄̅̅̅̅̅̅̅̅ = 𝑃⊙𝑄\par\vspace{2pt}\noindent
\textbf{(B)}~𝑃̅ $\oplus$𝑄= 𝑃⊙𝑄\par\vspace{2pt}\noindent
\textbf{(C)}~𝑃̅ $\oplus$𝑄̅ = 𝑃$\oplus$𝑄\par\vspace{2pt}\noindent
\textbf{(D)}~(𝑃$\oplus$𝑃̅) $\oplus$𝑄= (𝑃⊙𝑃̅) ⊙𝑄̅
\par\vspace{6pt}\hrule height 0.25pt\vspace{6pt}
\noindent\textbf{Q.5} \hfill {\footnotesize\textsf{[MCQ]}}
\par\vspace{2pt}
{\bfseries What is the missing number in the following sequence?
2, 12, 60, 240, 720, 1440, \_\_\_\_\_, 0 \par}
\par\vspace{2pt}\noindent
\begin{tabular*}{\columnwidth}{@{\extracolsep{\fill}}p{0.46\columnwidth}p{0.46\columnwidth}@{}}
\textbf{(A)}~2880 & \textbf{(B)}~1440 \\
\addlinespace[2pt]
\textbf{(C)}~720 & \textbf{(D)}~0
\end{tabular*}
\par\vspace{6pt}\hrule height 0.25pt\vspace{6pt}
\noindent\textbf{Q.5} \hfill {\footnotesize\textsf{[MCQ]}}
\par\vspace{2pt}
{\bfseries Consider the following processor design characteristics.
I.
Register-to-register arithmetic operations only
II.
Fixed-length instruction format
III.
Hardwired control unit
\par\smallskip
Which of the characteristics above are used in the design of a RISC processor? \par}
\par\vspace{2pt}\noindent
\begin{tabular*}{\columnwidth}{@{\extracolsep{\fill}}p{0.46\columnwidth}p{0.46\columnwidth}@{}}
\textbf{(A)}~I and II only & \textbf{(B)}~II and III only \\
\addlinespace[2pt]
\textbf{(C)}~I and III only & \textbf{(D)}~I, II and III
\end{tabular*}
\par\vspace{6pt}\hrule height 0.25pt\vspace{6pt}
\noindent\textbf{Q.6} \hfill {\footnotesize\textsf{[MCQ]}}
\par\vspace{2pt}
{\bfseries Let N be an NFA with n states. Let k be the number of states of a minimal DFA which is
equivalent to N. Which one of the following is necessarily true? \par}
\par\vspace{2pt}\noindent
\begin{tabular*}{\columnwidth}{@{\extracolsep{\fill}}p{0.46\columnwidth}p{0.46\columnwidth}@{}}
\textbf{(A)}~$k$$\ge$2$n$ & \textbf{(B)}~$k$$\ge$$n$ \\
\addlinespace[2pt]
\textbf{(C)}~$k$$\le$$n$2 & \textbf{(D)}~$k$$\le$2$n$
\end{tabular*}
\par\vspace{6pt}\hrule height 0.25pt\vspace{6pt}

\vspace{6pt}
\noindent{\footnotesize\textit{[Section Block: 2-mark questions]}}
\vspace{3pt}

\noindent\textbf{Q.6} \hfill {\footnotesize\textsf{[MCQ]}}
\par\vspace{2pt}
{\bfseries In appreciation of the social improvements completed in a town, a wealthy philanthropist
decided to gift Rs 750 to each male senior citizen in the town and Rs 1000 to each female
senior citizen. Altogether, there were 300 senior citizens eligible for this gift. However,
only 8/9th of the eligible men and 2/3rd of the eligible women claimed the gift. How much
money (in Rupees) did the philanthropist give away in total? \par}
\par\vspace{2pt}\noindent
\begin{tabular*}{\columnwidth}{@{\extracolsep{\fill}}p{0.46\columnwidth}p{0.46\columnwidth}@{}}
\textbf{(A)}~1,50,000 & \textbf{(B)}~2,00,000 \\
\addlinespace[2pt]
\textbf{(C)}~1,75,000 & \textbf{(D)}~1,51,000
\end{tabular*}
\par\vspace{6pt}\hrule height 0.25pt\vspace{6pt}

\vspace{6pt}
\noindent{\footnotesize\textit{[Section Block: 1-mark questions]}}
\vspace{3pt}

\noindent\textbf{Q.7} \hfill {\footnotesize\textsf{[MCQ]}}
\par\vspace{2pt}
{\bfseries The set of all recursively enumerable languages is \par}
\par\vspace{2pt}\noindent
\textbf{(A)}~closed under complementation.\par\vspace{2pt}\noindent
\textbf{(B)}~closed under intersection.\par\vspace{2pt}\noindent
\textbf{(C)}~a subset of the set of all recursive languages.\par\vspace{2pt}\noindent
\textbf{(D)}~an uncountable set.
\par\vspace{6pt}\hrule height 0.25pt\vspace{6pt}

\vspace{6pt}
\noindent{\footnotesize\textit{[Section Block: 2-mark questions]}}
\vspace{3pt}

\noindent\textbf{Q.7} \hfill {\footnotesize\textsf{[MCQ]}}
\par\vspace{2pt}
{\bfseries If 𝑝𝑞$r$$\ne$0 and 𝑝−$x$=
1
𝑞, 𝑞−$y$=
1
$r$, $r$−$z$=
1
𝑝 , what is the value of the product $x$$y$$z$? \par}
\par\vspace{2pt}\noindent
\textbf{(A)}~−1\par\vspace{2pt}\noindent
\textbf{(B)}~1 \newline 𝑝𝑞$r$\par\vspace{2pt}\noindent
\textbf{(C)}~1\par\vspace{2pt}\noindent
\textbf{(D)}~𝑝𝑞$r$
\par\vspace{6pt}\hrule height 0.25pt\vspace{6pt}

\vspace{6pt}
\noindent{\footnotesize\textit{[Section Block: 1-mark questions]}}
\vspace{3pt}

\noindent\textbf{Q.8} \hfill {\footnotesize\textsf{[MCQ]}}
\par\vspace{2pt}
{\bfseries Which one of the following statements is FALSE? \par}
\par\vspace{2pt}\noindent
\textbf{(A)}~Context-free grammar can be used to specify both lexical and syntax rules.\par\vspace{2pt}\noindent
\textbf{(B)}~Type checking is done before parsing.\par\vspace{2pt}\noindent
\textbf{(C)}~High-level language programs can be translated to different Intermediate \newline Representations.\par\vspace{2pt}\noindent
\textbf{(D)}~Arguments to a function can be passed using the program stack.
\par\vspace{6pt}\hrule height 0.25pt\vspace{6pt}

\vspace{6pt}
\noindent{\footnotesize\textit{[Section Block: 2-mark questions]}}
\vspace{3pt}

\noindent\textbf{Q.8} \hfill {\footnotesize\textsf{[MCQ]}}
\par\vspace{2pt}
{\bfseries In a party, 60\% of the invited guests are male and 40\% are female. If 80\% of the invited
guests attended the party and if all the invited female guests attended, what would be the
ratio of males to females among the attendees in the party? \par}
\par\vspace{2pt}\noindent
\begin{tabular*}{\columnwidth}{@{\extracolsep{\fill}}p{0.46\columnwidth}p{0.46\columnwidth}@{}}
\textbf{(A)}~2:3 & \textbf{(B)}~1:1 \\
\addlinespace[2pt]
\textbf{(C)}~3:2 & \textbf{(D)}~2:1
\end{tabular*}
\par\vspace{6pt}\hrule height 0.25pt\vspace{6pt}

\vspace{6pt}
\noindent{\footnotesize\textit{[Section Block: 1-mark questions]}}
\vspace{3pt}

\noindent\textbf{Q.9} \hfill {\footnotesize\textsf{[MCQ]}}
\par\vspace{2pt}
{\bfseries The following are some events that occur after a device controller issues an interrupt while
process L is under execution.
(P) The processor pushes the process status of L onto the control stack.
(Q) The processor finishes the execution of the current instruction.
(R) The processor executes the interrupt service routine.
(S) The processor pops the process status of L from the control stack.
(T) The processor loads the new PC value based on the interrupt.
Which one of the following is the correct order in which the events above occur? \par}
\par\vspace{2pt}\noindent
\begin{tabular*}{\columnwidth}{@{\extracolsep{\fill}}p{0.46\columnwidth}p{0.46\columnwidth}@{}}
\textbf{(A)}~QPTRS & \textbf{(B)}~PTRSQ \\
\addlinespace[2pt]
\textbf{(C)}~TRPQS & \textbf{(D)}~QTPRS
\end{tabular*}
\par\vspace{6pt}\hrule height 0.25pt\vspace{6pt}
\noindent\textbf{Q.10} \hfill {\footnotesize\textsf{[MCQ]}}
\par\vspace{2pt}
{\bfseries Consider a process executing on an operating system that uses demand paging. The average
time for a memory access in the system is M units if the corresponding memory page is
available in memory, and D units if the memory access causes a page fault. It has been
experimentally measured that the average time taken for a memory access in the process is
X units.
\par\smallskip
Which one of the following is the correct expression for the page fault rate experienced by
the process? \par}
\par\vspace{2pt}\noindent
\begin{tabular*}{\columnwidth}{@{\extracolsep{\fill}}p{0.46\columnwidth}p{0.46\columnwidth}@{}}
\textbf{(A)}~(D -- M) / (X -- M) & \textbf{(B)}~(X -- M) / (D -- M) \\
\addlinespace[2pt]
\textbf{(C)}~(D -- X) / (D -- M) & \textbf{(D)}~(X -- M) / (D -- X)
\end{tabular*}
\par\vspace{6pt}\hrule height 0.25pt\vspace{6pt}

\vspace{6pt}
\noindent{\footnotesize\textit{[Section Block: 2-mark questions]}}
\vspace{3pt}

\noindent\textbf{Q.10} \hfill {\footnotesize\textsf{[MCQ]}}
\par\vspace{2pt}
{\bfseries A six sided unbiased die with four green faces and two red faces is rolled seven times.
Which of the following combinations is the most likely outcome of the experiment? \par}
\par\vspace{2pt}\noindent
\textbf{(A)}~Three green faces and four red faces.\par\vspace{2pt}\noindent
\textbf{(B)}~Four green faces and three red faces.\par\vspace{2pt}\noindent
\textbf{(C)}~Five green faces and two red faces.\par\vspace{2pt}\noindent
\textbf{(D)}~Six green faces and one red face.
\par\vspace{6pt}\hrule height 0.25pt\vspace{6pt}

\vspace{6pt}
\noindent{\footnotesize\textit{[Section Block: 1-mark questions]}}
\vspace{3pt}

\noindent\textbf{Q.11} \hfill {\footnotesize\textsf{[MCQ]}}
\par\vspace{2pt}
{\bfseries In an Entity-Relationship (ER) model, suppose 𝑅 is a many-to-one relationship from entity
set E1 to entity set E2. Assume that E1 and E2 participate totally in 𝑅 and that the cardinality
of E1 is greater than the cardinality of E2.
\par\smallskip
Which one of the following is true about 𝑅? \par}
\par\vspace{2pt}\noindent
\textbf{(A)}~Every entity in E1 is associated with exactly one entity in E2.\par\vspace{2pt}\noindent
\textbf{(B)}~Some entity in E1 is associated with more than one entity in E2.\par\vspace{2pt}\noindent
\textbf{(C)}~Every entity in E2 is associated with exactly one entity in E1.\par\vspace{2pt}\noindent
\textbf{(D)}~Every entity in E2 is associated with at most one entity in E1.
\par\vspace{6pt}\hrule height 0.25pt\vspace{6pt}
\noindent\textbf{Q.12} \hfill {\footnotesize\textsf{[MCQ]}}
\par\vspace{2pt}
{\bfseries Consider the following two tables and four queries in SQL.
Book (isbn, bname), Stock (isbn, copies)
\par\smallskip
Query 1: SELECT B.isbn, S.copies
FROM Book B INNER JOIN Stock S
ON B.isbn = S.isbn;
\par\smallskip
Query 2: SELECT B.isbn, S.copies
FROM Book B LEFT OUTER JOIN Stock S
ON B.isbn = S.isbn;
\par\smallskip
Query 3: SELECT B.isbn, S.copies
FROM Book B RIGHT OUTER JOIN Stock S
ON B.isbn = S.isbn;
\par\smallskip
Query 4: SELECT B.isbn, S.copies
FROM Book B FULL OUTER JOIN Stock S
ON B.isbn = S.isbn;
\par\smallskip
Which one of the queries above is certain to have an output that is a superset of the outputs
of the other three queries? \par}
\par\vspace{2pt}\noindent
\begin{tabular*}{\columnwidth}{@{\extracolsep{\fill}}p{0.46\columnwidth}p{0.46\columnwidth}@{}}
\textbf{(A)}~Query 1 & \textbf{(B)}~Query 2 \\
\addlinespace[2pt]
\textbf{(C)}~Query 3 & \textbf{(D)}~Query 4
\end{tabular*}
\par\vspace{6pt}\hrule height 0.25pt\vspace{6pt}
\noindent\textbf{Q.13} \hfill {\footnotesize\textsf{[MCQ]}}
\par\vspace{2pt}
{\bfseries Match the following:
Field
Length in bits
P. UDP Header's Port Number
I. 48
Q. Ethernet MAC Address
II. 8
R. IPv6 Next Header
III. 32
S. TCP Header's Sequence Number
IV. 16 \par}
\par\vspace{2pt}\noindent
\textbf{(A)}~P-III, Q-IV, R-II, S-I\par\vspace{2pt}\noindent
\textbf{(B)}~P-II, Q-I, R-IV, S-III\par\vspace{2pt}\noindent
\textbf{(C)}~P-IV, Q-I, R-II, S-III\par\vspace{2pt}\noindent
\textbf{(D)}~P-IV, Q-I, R-III, S-II
\par\vspace{6pt}\hrule height 0.25pt\vspace{6pt}
\noindent\textbf{Q.14} \hfill {\footnotesize\textsf{[MCQ]}}
\par\vspace{2pt}
{\bfseries Consider the following statements regarding the slow start phase of the TCP congestion
control algorithm. Note that cwnd stands for the TCP congestion window and MSS denotes
the Maximum Segment Size.
(i) The cwnd increases by 2 MSS on every successful acknowledgment.
(ii) The cwnd approximately doubles on every successful acknowledgement.
(iii) The cwnd increases by 1 MSS every round trip time.
(iv) The cwnd approximately doubles every round trip time.
\par\smallskip
Which one of the following is correct? \par}
\par\vspace{2pt}\noindent
\textbf{(A)}~Only (ii) and (iii) are true\par\vspace{2pt}\noindent
\textbf{(B)}~Only (i) and (iii) are true\par\vspace{2pt}\noindent
\textbf{(C)}~Only (iv) is true\par\vspace{2pt}\noindent
\textbf{(D)}~Only (i) and (iv) are true
\par\vspace{6pt}\hrule height 0.25pt\vspace{6pt}
\noindent\textbf{Q.15} \hfill {\footnotesize\textsf{[NAT]}}
\par\vspace{2pt}
{\bfseries Two people, P and Q, decide to independently roll two identical dice, each with 6 faces,
numbered 1 to 6. The person with the lower number wins. In case of a tie, they roll the dice
repeatedly until there is no tie. Define a trial as a throw of the dice by P and Q. Assume that
all 6 numbers on each dice are equi-probable and that all trials are independent. The
probability (rounded to 3 decimal places) that one of them wins on the third trial is \_\_\_\_\_. \par}
\par\vspace{3pt}\noindent\textbf{Answer:} \rule{3.2cm}{0.4pt}
\par\vspace{6pt}\hrule height 0.25pt\vspace{6pt}
\noindent\textbf{Q.16} \hfill {\footnotesize\textsf{[NAT]}}
\par\vspace{2pt}
{\bfseries The value of ∫
$x$cos($x$2)𝑑$x$
𝜋/4
0
correct to three decimal places
(assuming that 𝜋= 3.14 ) is \_\_\_\_\_. \par}
\par\vspace{3pt}\noindent\textbf{Answer:} \rule{3.2cm}{0.4pt}
\par\vspace{6pt}\hrule height 0.25pt\vspace{6pt}
\noindent\textbf{Q.17} \hfill {\footnotesize\textsf{[NAT]}}
\par\vspace{2pt}
{\bfseries Consider a matrix
T
A
uv
=
where
1
1
,
2
1
u
v






=
=
. Note that
Tv denotes the
transpose of v. The largest eigenvalue of A is \_\_\_\_\_. \par}
\par\vspace{3pt}\noindent\textbf{Answer:} \rule{3.2cm}{0.4pt}
\par\vspace{6pt}\hrule height 0.25pt\vspace{6pt}
\noindent\textbf{Q.18} \hfill {\footnotesize\textsf{[NAT]}}
\par\vspace{2pt}
{\bfseries The chromatic number of the following graph is \_\_\_\_\_\_\_. \par}
\par\vspace{3pt}\noindent\textbf{Answer:} \rule{3.2cm}{0.4pt}
\par\vspace{6pt}\hrule height 0.25pt\vspace{6pt}
\noindent\textbf{Q.19} \hfill {\footnotesize\textsf{[NAT]}}
\par\vspace{2pt}
{\bfseries Let G be a finite group on 84 elements. The size of a largest possible proper subgroup of G
is \_\_\_\_\_\_\_\_. \par}
\par\vspace{3pt}\noindent\textbf{Answer:} \rule{3.2cm}{0.4pt}
\par\vspace{6pt}\hrule height 0.25pt\vspace{6pt}
\noindent\textbf{Q.20} \hfill {\footnotesize\textsf{[NAT]}}
\par\vspace{2pt}
{\bfseries The postorder traversal of a binary tree is 8,9,6,7,4,5,2,3,1. The inorder traversal of the same
tree is 8,6,9,4,7,2,5,1,3. The height of a tree is the length of the longest path from the root to
any leaf. The height of the binary tree above is \_\_\_\_\_\_. \par}
\par\vspace{3pt}\noindent\textbf{Answer:} \rule{3.2cm}{0.4pt}
\par\vspace{6pt}\hrule height 0.25pt\vspace{6pt}
\noindent\textbf{Q.21} \hfill {\footnotesize\textsf{[NAT]}}
\par\vspace{2pt}
{\bfseries Consider the following C program:
\#include <stdio.h>
\par\smallskip
int counter = 0;
\par\smallskip
int calc (int a, int b) \{
int c;
\par\smallskip
counter++;
if (b==3) return (a*a*a);
else \{
c = calc(a, b/3);
return (c*c*c);
\}
\}
\par\smallskip
int main ()\{
calc(4, 81);
printf ("\%d", counter);
\}
\par\smallskip
The output of this program is \_\_\_\_\_. \par}
\par\vspace{3pt}\noindent\textbf{Answer:} \rule{3.2cm}{0.4pt}
\par\vspace{6pt}\hrule height 0.25pt\vspace{6pt}
\noindent\textbf{Q.23} \hfill {\footnotesize\textsf{[NAT]}}
\par\vspace{2pt}
{\bfseries A 32-bit wide main memory unit with a capacity of 1 GB is built using 256M $\times$ 4-bit DRAM
chips. The number of rows of memory cells in the DRAM chip is 214. The time taken to
perform one refresh operation is 50 nanoseconds. The refresh period is 2 milliseconds. The
percentage (rounded to the closest integer) of the time available for performing the memory
read/write operations in the main memory unit is \_\_\_\_\_\_\_\_\_\_. \par}
\par\vspace{3pt}\noindent\textbf{Answer:} \rule{3.2cm}{0.4pt}
\par\vspace{6pt}\hrule height 0.25pt\vspace{6pt}
\noindent\textbf{Q.24} \hfill {\footnotesize\textsf{[NAT]}}
\par\vspace{2pt}
{\bfseries Consider a system with 3 processes that share 4 instances of the same resource type. Each
process can request a maximum of 𝐾 instances. Resource instances can be requested and
released only one at a time. The largest value of 𝐾 that will always avoid deadlock is \_\_\_\_. \par}
\par\vspace{3pt}\noindent\textbf{Answer:} \rule{3.2cm}{0.4pt}
\par\vspace{6pt}\hrule height 0.25pt\vspace{6pt}
\noindent\textbf{Q.25} \hfill {\footnotesize\textsf{[NAT]}}
\par\vspace{2pt}
{\bfseries Consider a long-lived TCP session with an end-to-end bandwidth of 1 Gbps (= 109 bits-per-
second). The session starts with a sequence number of 1234. The minimum time (in seconds,
rounded to the closest integer) before this sequence number can be used again is \_\_\_\_\_\_\_. \par}
\par\vspace{3pt}\noindent\textbf{Answer:} \rule{3.2cm}{0.4pt}
\par\vspace{6pt}\hrule height 0.25pt\vspace{6pt}

\vspace{6pt}
\noindent{\footnotesize\textit{[Section Block: 2-mark questions]}}
\vspace{3pt}

\noindent\textbf{Q.26} \hfill {\footnotesize\textsf{[MCQ]}}
\par\vspace{2pt}
{\bfseries Consider a matrix P whose only eigenvectors are the multiples of 1
4



.
Consider the following statements.
(I) P does not have an inverse
(II) P has a repeated eigenvalue
(III) P cannot be diagonalized
Which one of the following options is correct? \par}
\par\vspace{2pt}\noindent
\textbf{(A)}~Only I and III are necessarily true\par\vspace{2pt}\noindent
\textbf{(B)}~Only II is necessarily true\par\vspace{2pt}\noindent
\textbf{(C)}~Only I and II are necessarily true\par\vspace{2pt}\noindent
\textbf{(D)}~Only II and III are necessarily true
\par\vspace{6pt}\hrule height 0.25pt\vspace{6pt}
\noindent\textbf{Q.27} \hfill {\footnotesize\textsf{[MCQ]}}
\par\vspace{2pt}
{\bfseries Let N be the set of natural numbers. Consider the following sets.
P: Set of Rational numbers (positive and negative)
Q: Set of functions from \{0, 1\} to N
R: Set of functions from N to \{0, 1\}
S: Set of finite subsets of N.
\par\smallskip
Which of the sets above are countable? \par}
\par\vspace{2pt}\noindent
\begin{tabular*}{\columnwidth}{@{\extracolsep{\fill}}p{0.46\columnwidth}p{0.46\columnwidth}@{}}
\textbf{(A)}~Q and S only & \textbf{(B)}~P and S only \\
\addlinespace[2pt]
\textbf{(C)}~P and R only & \textbf{(D)}~P, Q and S only
\end{tabular*}
\par\vspace{6pt}\hrule height 0.25pt\vspace{6pt}
\noindent\textbf{Q.28} \hfill {\footnotesize\textsf{[MCQ]}}
\par\vspace{2pt}
{\bfseries Consider the first-order logic sentence
𝜑$\equiv$ ∃𝑠∃𝑡∃𝑢∀𝑣∀𝑤∀$x$∀$y$ 𝜓(𝑠, 𝑡, 𝑢, 𝑣, 𝑤, $x$, $y$)
\par\smallskip
where 𝜓(𝑠, 𝑡, 𝑢, 𝑣, 𝑤, $x$, $y$) is a quantifier-free first-order logic formula using only predicate
symbols, and possibly equality, but no function symbols. Suppose 𝜑 has a model with a
universe containing 7 elements.
\par\smallskip
Which one of the following statements is necessarily true? \par}
\par\vspace{2pt}\noindent
\textbf{(A)}~There exists at least one model of 𝜑 with universe of size less than or equal to 3.\par\vspace{2pt}\noindent
\textbf{(B)}~There exists no model of 𝜑 with universe of size less than or equal to 3.\par\vspace{2pt}\noindent
\textbf{(C)}~There exists no model of 𝜑 with universe of size greater than 7.\par\vspace{2pt}\noindent
\textbf{(D)}~Every model of 𝜑 has a universe of size equal to 7.
\par\vspace{6pt}\hrule height 0.25pt\vspace{6pt}
\noindent\textbf{Q.29} \hfill {\footnotesize\textsf{[MCQ]}}
\par\vspace{2pt}
{\bfseries Consider the following C program:
\#include<stdio.h>
\par\smallskip
void fun1(char *s1, char *s2)\{
\par\smallskip
char *tmp;
\par\smallskip
tmp = s1;
\par\smallskip
s1 = s2;
\par\smallskip
s2 = tmp;
\}
void fun2(char **s1, char **s2)\{
\par\smallskip
char *tmp;
\par\smallskip
tmp = *s1;
\par\smallskip
*s1 = *s2;
\par\smallskip
*s2 = tmp;
\}
int main()\{
\par\smallskip
char *str1 = "Hi", *str2 = "Bye";
\par\smallskip
fun1(str1, str2); printf("\%s \%s ", str1, str2);
\par\smallskip
fun2(\&str1, \&str2); printf("\%s \%s", str1, str2);
\par\smallskip
return 0;
\}
\par\smallskip
The output of the program above is \par}
\par\vspace{2pt}\noindent
\begin{tabular*}{\columnwidth}{@{\extracolsep{\fill}}p{0.46\columnwidth}p{0.46\columnwidth}@{}}
\textbf{(A)}~Hi Bye Bye Hi & \textbf{(B)}~Hi Bye Hi Bye \\
\addlinespace[2pt]
\textbf{(C)}~Bye Hi Hi Bye & \textbf{(D)}~Bye Hi Bye Hi
\end{tabular*}
\par\vspace{6pt}\hrule height 0.25pt\vspace{6pt}
\noindent\textbf{Q.30} \hfill {\footnotesize\textsf{[MCQ]}}
\par\vspace{2pt}
{\bfseries Let G be a simple undirected graph. Let TD be a depth first search tree of G. Let TB be a
breadth first search tree of G. Consider the following statements.
(I) No edge of G is a cross edge with respect to TD. (A cross edge in G is between two
nodes neither of which is an ancestor of the other in TD.)
(II) For every edge (u,v) of G, if u is at depth i and v is at depth j in TB, then |$i$−$j$| = 1.
Which of the statements above must necessarily be true? \par}
\par\vspace{2pt}\noindent
\begin{tabular*}{\columnwidth}{@{\extracolsep{\fill}}p{0.46\columnwidth}p{0.46\columnwidth}@{}}
\textbf{(A)}~I only & \textbf{(B)}~II only \\
\addlinespace[2pt]
\textbf{(C)}~Both I and II & \textbf{(D)}~Neither I nor II
\end{tabular*}
\par\vspace{6pt}\hrule height 0.25pt\vspace{6pt}
\noindent\textbf{Q.31} \hfill {\footnotesize\textsf{[MCQ]}}
\par\vspace{2pt}
{\bfseries Assume that multiplying a matrix G1 of dimension 𝑝 $\times$ 𝑞 with another matrix G2 of
dimension 𝑞$\times$ $r$ requires 𝑝𝑞$r$ scalar multiplications. Computing the product of n matrices
G1G2G3…Gn can be done by parenthesizing in different ways. Define Gi Gi+1 as an explicitly
computed pair for a given paranthesization if they are directly multiplied. For example, in
the
matrix
multiplication
chain
G1G2G3G4G5G6
using
parenthesization
(G1(G2G3))(G4(G5G6)), G2G3 and G5G6 are the only explicitly computed pairs.
\par\smallskip
Consider a matrix multiplication chain F1F2F3F4F5, where matrices F1, F2, F3, F4 and F5 are
of dimensions 2$\times$25, 25$\times$3, 3$\times$16, 16$\times$1 and 1$\times$1000, respectively. In the parenthesization of
F1F2F3F4F5 that minimizes the total number of scalar multiplications, the explicitly
computed pairs is/are \par}
\par\vspace{2pt}\noindent
\begin{tabular*}{\columnwidth}{@{\extracolsep{\fill}}p{0.46\columnwidth}p{0.46\columnwidth}@{}}
\textbf{(A)}~F1F2 and F3F4 only & \textbf{(B)}~F2F3 only \\
\addlinespace[2pt]
\textbf{(C)}~F3F4 only & \textbf{(D)}~F1F2 and F4F5 only
\end{tabular*}
\par\vspace{6pt}\hrule height 0.25pt\vspace{6pt}
\noindent\textbf{Q.32} \hfill {\footnotesize\textsf{[MCQ]}}
\par\vspace{2pt}
{\bfseries Consider the following C code. Assume that unsigned long int type length is 64
bits.
unsigned long int fun(unsigned long int n)\{
unsigned long int i, j = 0, sum = 0;
for (i = n; i > 1; i = i/2) j++;
for ( ; j > 1; j = j/2) sum++;
return(sum);
\}
\par\smallskip
The value returned when we call fun with the input 240 is \par}
\par\vspace{2pt}\noindent
\begin{tabular*}{\columnwidth}{@{\extracolsep{\fill}}p{0.46\columnwidth}p{0.46\columnwidth}@{}}
\textbf{(A)}~4 & \textbf{(B)}~5 \\
\addlinespace[2pt]
\textbf{(C)}~6 & \textbf{(D)}~40
\end{tabular*}
\par\vspace{6pt}\hrule height 0.25pt\vspace{6pt}
\noindent\textbf{Q.33} \hfill {\footnotesize\textsf{[MCQ]}}
\par\vspace{2pt}
{\bfseries Consider the unsigned 8-bit fixed point binary number representation below,
b7 b6 b5 b4 b3 . b2 b1 b0
where the position of the binary point is between b3 and b2. Assume b7 is the most significant
bit. Some of the decimal numbers listed below cannot be represented exactly in the above
representation:
(i)
31.500 (ii) 0.875 (iii) 12.100 (iv) 3.001
\par\smallskip
Which one of the following statements is true? \par}
\par\vspace{2pt}\noindent
\textbf{(A)}~None of (i), (ii), (iii), (iv) can be exactly represented\par\vspace{2pt}\noindent
\textbf{(B)}~Only (ii) cannot be exactly represented\par\vspace{2pt}\noindent
\textbf{(C)}~Only (iii) and (iv) cannot be exactly represented\par\vspace{2pt}\noindent
\textbf{(D)}~Only (i) and (ii) cannot be exactly represented
\par\vspace{6pt}\hrule height 0.25pt\vspace{6pt}
\noindent\textbf{Q.34} \hfill {\footnotesize\textsf{[MCQ]}}
\par\vspace{2pt}
{\bfseries The size of the physical address space of a processor is 2𝑃 bytes. The word length is
2𝑊 bytes. The capacity of cache memory is 2𝑁 bytes. The size of each cache block is
2𝑀 words. For a 𝐾-way set-associative cache memory, the length (in number of bits) of the
tag field is \par}
\par\vspace{2pt}\noindent
\begin{tabular*}{\columnwidth}{@{\extracolsep{\fill}}p{0.46\columnwidth}p{0.46\columnwidth}@{}}
\textbf{(A)}~𝑃−𝑁− log2 𝐾 & \textbf{(B)}~𝑃−𝑁+ log2 𝐾 \\
\addlinespace[2pt]
\textbf{(C)}~𝑃−𝑁−𝑀−𝑊− log2 𝐾 & \textbf{(D)}~𝑃−𝑁−𝑀−𝑊+ log2 𝐾
\end{tabular*}
\par\vspace{6pt}\hrule height 0.25pt\vspace{6pt}
\noindent\textbf{Q.35} \hfill {\footnotesize\textsf{[MCQ]}}
\par\vspace{2pt}
{\bfseries Consider the following languages:
I. \{𝑎𝑚𝑏$n$𝑐𝑝𝑑𝑞| 𝑚+ 𝑝= $n$+ 𝑞, where 𝑚, $n$, 𝑝, 𝑞$\ge$0\}
II. \{𝑎𝑚𝑏$n$𝑐𝑝𝑑𝑞| 𝑚= $n$ and 𝑝= 𝑞, where 𝑚, $n$, 𝑝, 𝑞$\ge$0\}
III. \{𝑎𝑚𝑏$n$𝑐𝑝𝑑𝑞| 𝑚= $n$= 𝑝 and 𝑝$\ne$𝑞, where 𝑚, $n$, 𝑝, 𝑞$\ge$0\}
IV. \{𝑎𝑚𝑏$n$𝑐𝑝𝑑𝑞| 𝑚$n$= 𝑝+ 𝑞, where 𝑚, $n$, 𝑝, 𝑞$\ge$0\}
Which of the languages above are context-free? \par}
\par\vspace{2pt}\noindent
\begin{tabular*}{\columnwidth}{@{\extracolsep{\fill}}p{0.46\columnwidth}p{0.46\columnwidth}@{}}
\textbf{(A)}~I and IV only & \textbf{(B)}~I and II only \\
\addlinespace[2pt]
\textbf{(C)}~II and III only & \textbf{(D)}~II and IV only
\end{tabular*}
\par\vspace{6pt}\hrule height 0.25pt\vspace{6pt}
\noindent\textbf{Q.36} \hfill {\footnotesize\textsf{[MCQ]}}
\par\vspace{2pt}
{\bfseries Consider the following problems. 𝐿(𝐺) denotes the language generated by a grammar 𝐺.
𝐿(𝑀) denotes the language accepted by a machine 𝑀.
(I)
For an unrestricted grammar 𝐺 and a string 𝑤, whether 𝑤$\in$𝐿(𝐺)
(II)
Given a Turing machine M, whether L(M) is regular
(III)
Given two grammars 𝐺1 and 𝐺2, whether 𝐿(𝐺1) = 𝐿(𝐺2)
(IV)
Given an NFA N, whether there is a deterministic PDA P such that N and P
accept the same language.
\par\smallskip
Which one of the following statements is correct? \par}
\par\vspace{2pt}\noindent
\textbf{(A)}~Only I and II are undecidable\par\vspace{2pt}\noindent
\textbf{(B)}~Only III is undecidable\par\vspace{2pt}\noindent
\textbf{(C)}~Only II and IV are undecidable\par\vspace{2pt}\noindent
\textbf{(D)}~Only I, II and III are undecidable
\par\vspace{6pt}\hrule height 0.25pt\vspace{6pt}
\noindent\textbf{Q.37} \hfill {\footnotesize\textsf{[MCQ]}}
\par\vspace{2pt}
{\bfseries A lexical analyzer uses the following patterns to recognize three tokens T1, T2, and T3 over
the alphabet \{a,b,c\}.
\par\smallskip
𝑇1: 𝑎? (𝑏|𝑐)∗𝑎
𝑇2: 𝑏? (𝑎|𝑐)∗𝑏
𝑇3: 𝑐? (𝑏|𝑎)∗𝑐
\par\smallskip
Note that `x?' means 0 or 1 occurrence of the symbol x. Note also that the analyzer outputs
the token that matches the longest possible prefix.
\par\smallskip
If the string 𝑏𝑏𝑎𝑎𝑐𝑎𝑏𝑐 is processed by the analyzer, which one of the following is the
sequence of tokens it outputs? \par}
\par\vspace{2pt}\noindent
\begin{tabular*}{\columnwidth}{@{\extracolsep{\fill}}p{0.46\columnwidth}p{0.46\columnwidth}@{}}
\textbf{(A)}~𝑇1𝑇2𝑇3 & \textbf{(B)}~𝑇1𝑇1𝑇3 \\
\addlinespace[2pt]
\textbf{(C)}~𝑇2𝑇1𝑇3 & \textbf{(D)}~𝑇3𝑇3
\end{tabular*}
\par\vspace{6pt}\hrule height 0.25pt\vspace{6pt}
\noindent\textbf{Q.39} \hfill {\footnotesize\textsf{[MCQ]}}
\par\vspace{2pt}
{\bfseries In a system, there are three types of resources: E, F and G. Four processes P0, P1, P2 and P3
execute concurrently. At the outset, the processes have declared their maximum resource
requirements using a matrix named Max as given below. For example, Max[P2,F] is the
maximum number of instances of F that P2 would require. The number of instances of the
resources allocated to the various processes at any given state is given by a matrix named
Allocation.
\par\smallskip
Consider a state of the system with the Allocation matrix as shown below, and in which 3
instances of E and 3 instances of F are the only resources available.
Allocation
\par\smallskip
Max
\par\smallskip
E F G
\par\smallskip
E
F G
P0
1 0
1
P0
4
3
1
P1 1 1
2
P1
2
1
4
P2 1 0
3
P2
1
3
3
P3 2 0
0
P3
5
4
1
\par\smallskip
From the perspective of deadlock avoidance, which one of the following is true? \par}
\par\vspace{2pt}\noindent
\textbf{(A)}~The system is in safe state.\par\vspace{2pt}\noindent
\textbf{(B)}~The system is not in safe state, but would be safe if one more instance of E were \newline available\par\vspace{2pt}\noindent
\textbf{(C)}~The system is not in safe state, but would be safe if one more instance of F were \newline available\par\vspace{2pt}\noindent
\textbf{(D)}~The system is not in safe state, but would be safe if one more instance of G were \newline available
\par\vspace{6pt}\hrule height 0.25pt\vspace{6pt}
\noindent\textbf{Q.40} \hfill {\footnotesize\textsf{[MCQ]}}
\par\vspace{2pt}
{\bfseries Consider the following solution to the producer-consumer synchronization problem. The
shared buffer size is 𝑁. Three semaphores empty, full and mutex are defined with respective
initial values of 0, 𝑁 and 1. Semaphore empty denotes the number of available slots in the
buffer, for the consumer to read from. Semaphore full denotes the number of available slots
in the buffer, for the producer to write to. The placeholder variables, denoted by P, Q, R,
and S, in the code below can be assigned either empty or full. The valid semaphore
operations are: wait() and signal().
Producer:
Consumer:
\par\smallskip
do\{
wait(P);
wait(mutex);
//Add item to buffer
signal(mutex);
signal(Q);
\}while(1);
do\{
wait(R);
wait(mutex);
//Consume item from buffer
signal(mutex);
signal(S);
\}while(1);
\par\smallskip
Which one of the following assignments to P, Q, R and S will yield the correct solution? \par}
\par\vspace{2pt}\noindent
\textbf{(A)}~P: full, Q: full, R: empty, S: empty\par\vspace{2pt}\noindent
\textbf{(B)}~P: empty, Q: empty, R: full, S: full\par\vspace{2pt}\noindent
\textbf{(C)}~P: full, Q: empty, R: empty, S: full\par\vspace{2pt}\noindent
\textbf{(D)}~P: empty, Q: full, R: full, S: empty
\par\vspace{6pt}\hrule height 0.25pt\vspace{6pt}
\noindent\textbf{Q.41} \hfill {\footnotesize\textsf{[MCQ]}}
\par\vspace{2pt}
{\bfseries Consider the relations r(A, B) and s(B, C), where s.B is a primary key and r.B is a foreign
key referencing s.B. Consider the query
\par\smallskip
Q: $r$⋈(𝜎𝐵<5(𝑠))
\par\smallskip
Let LOJ denote the natural left outer-join operation. Assume that r and s contain no null
values.
\par\smallskip
Which one of the following queries is NOT equivalent to Q? \par}
\par\vspace{2pt}\noindent
\begin{tabular*}{\columnwidth}{@{\extracolsep{\fill}}p{0.46\columnwidth}p{0.46\columnwidth}@{}}
\textbf{(A)}~𝜎𝐵<5($r$⋈𝑠) & \textbf{(B)}~𝜎𝐵<5($r$ 𝐿𝑂𝐽 𝑠) \\
\addlinespace[2pt]
\textbf{(C)}~$r$ 𝐿𝑂𝐽 (𝜎𝐵<5(𝑠)) & \textbf{(D)}~𝜎𝐵<5($r$) 𝐿𝑂𝐽 𝑠
\end{tabular*}
\par\vspace{6pt}\hrule height 0.25pt\vspace{6pt}
\noindent\textbf{Q.42} \hfill {\footnotesize\textsf{[MCQ]}}
\par\vspace{2pt}
{\bfseries Consider the following four relational schemas. For each schema, all non-trivial functional
dependencies are listed. The underlined attributes are the respective primary keys.
\par\smallskip
Schema I: Registration (rollno, courses)
Field `courses' is a set-valued attribute containing the set of courses a student has
registered for.
Non-trivial functional dependency:
rollno $\to$ courses
\par\smallskip
Schema II: Registration (rollno, courseid, email)
Non-trivial functional dependencies:
rollno, courseid $\to$ email
email $\to$ rollno
\par\smallskip
Schema III: Registration (rollno, courseid, marks, grade)
Non-trivial functional dependencies:
rollno, courseid $\to$ marks, grade
marks $\to$ grade
\par\smallskip
Schema IV: Registration (rollno, courseid, credit)
Non-trivial functional dependencies:
rollno, courseid $\to$ credit
courseid $\to$ credit
\par\smallskip
Which one of the relational schemas above is in 3NF but not in BCNF? \par}
\par\vspace{2pt}\noindent
\begin{tabular*}{\columnwidth}{@{\extracolsep{\fill}}p{0.46\columnwidth}p{0.46\columnwidth}@{}}
\textbf{(A)}~Schema I & \textbf{(B)}~Schema II \\
\addlinespace[2pt]
\textbf{(C)}~Schema III & \textbf{(D)}~Schema IV
\end{tabular*}
\par\vspace{6pt}\hrule height 0.25pt\vspace{6pt}
\noindent\textbf{Q.43} \hfill {\footnotesize\textsf{[NAT]}}
\par\vspace{2pt}
{\bfseries Let G be a graph with 100! vertices, with each vertex labelled by a distinct permutation of
the numbers 1,2, … , 100. There is an edge between vertices 𝑢 and 𝑣 if and only if the label
of 𝑢 can be obtained by swapping two adjacent numbers in the label of 𝑣. Let $y$ denote the
degree of a vertex in G, and $z$ denote the number of connected components in G.
Then, $y$ + 10$z$ = \_\_\_\_\_. \par}
\par\vspace{3pt}\noindent\textbf{Answer:} \rule{3.2cm}{0.4pt}
\par\vspace{6pt}\hrule height 0.25pt\vspace{6pt}
\noindent\textbf{Q.44} \hfill {\footnotesize\textsf{[NAT]}}
\par\vspace{2pt}
{\bfseries Consider Guwahati (G) and Delhi (D) whose temperatures can be classified as high (𝐻),
medium (𝑀) and low (𝐿). Let 𝑃(𝐻𝐺) denote the probability that Guwahati has high
temperature. Similarly, 𝑃(𝑀𝐺) and 𝑃(𝐿𝐺) denotes the probability of Guwahati having
medium and low temperatures respectively. Similarly, we use 𝑃(𝐻𝐷), 𝑃(𝑀𝐷) and 𝑃(𝐿𝐷) for
Delhi.
\par\smallskip
The following table gives the conditional probabilities for Delhi's temperature given
Guwahati's temperature.
𝐻𝐷
𝑀𝐷
𝐿𝐷
𝐻𝐺
0.40
0.48
0.12
𝑀𝐺
0.10
0.65
0.25
𝐿𝐺
0.01
0.50
0.49
\par\smallskip
Consider the first row in the table above. The first entry denotes that if Guwahati has high
temperature (𝐻𝐺) then the probability of Delhi also having a high temperature (𝐻𝐷) is 0.40;
i.e., 𝑃(𝐻𝐷|𝐻𝐺) = 0.40. Similarly, the next two entries are 𝑃(𝑀𝐷|𝐻𝐺) = 0.48 and
𝑃(𝐿𝐷|𝐻𝐺) = 0.12. Similarly for the other rows.
\par\smallskip
If it is known that 𝑃(𝐻𝐺) = 0.2, 𝑃(𝑀𝐺) = 0.5, and 𝑃(𝐿𝐺) = 0.3, then the probability
(correct to two decimal places) that Guwahati has high temperature given that Delhi has high
temperature is \_\_\_\_\_\_\_. \par}
\par\vspace{3pt}\noindent\textbf{Answer:} \rule{3.2cm}{0.4pt}
\par\vspace{6pt}\hrule height 0.25pt\vspace{6pt}
\noindent\textbf{Q.45} \hfill {\footnotesize\textsf{[NAT]}}
\par\vspace{2pt}
{\bfseries Consider the following program written in pseudo-code. Assume that x and y are integers.
Count(x,y) \{
if (y $\ne$ 1)\{
if (x $\ne$ 1) \{
print("*");
Count(x/2, y);
\}
else \{
y = y-1;
Count(1024, y);
\}
\}
\}
\par\smallskip
The number of times that the print statement is executed by the call
Count(1024,1024) is \_\_\_\_\_. \par}
\par\vspace{3pt}\noindent\textbf{Answer:} \rule{3.2cm}{0.4pt}
\par\vspace{6pt}\hrule height 0.25pt\vspace{6pt}
\noindent\textbf{Q.46} \hfill {\footnotesize\textsf{[NAT]}}
\par\vspace{2pt}
{\bfseries The number of possible min-heaps containing each value from \{1, 2, 3, 4, 5, 6, 7\}exactly once
is \_\_\_\_\_. \par}
\par\vspace{3pt}\noindent\textbf{Answer:} \rule{3.2cm}{0.4pt}
\par\vspace{6pt}\hrule height 0.25pt\vspace{6pt}
\noindent\textbf{Q.48} \hfill {\footnotesize\textsf{[NAT]}}
\par\vspace{2pt}
{\bfseries Consider the weights and values of items listed below. Note that there is only one unit of
each item.
Item number
Weight
(in Kgs)
Value
(in Rupees)
1
10
60
2
7
28
3
4
20
4
2
24
\par\smallskip
The task is to pick a subset of these items such that their total weight is no more than 11 Kgs
and their total value is maximized. Moreover, no item may be split. The total value of items
picked by an optimal algorithm is denoted by 𝑉opt. A greedy algorithm sorts the items by
their value-to-weight ratios in descending order and packs them greedily, starting from the
first item in the ordered list. The total value of items picked by the greedy algorithm is
denoted by 𝑉greedy.
\par\smallskip
The value of 𝑉opt −𝑉greedy is \_\_\_\_\_\_\_\_\_\_\_\_. \par}
\par\vspace{3pt}\noindent\textbf{Answer:} \rule{3.2cm}{0.4pt}
\par\vspace{6pt}\hrule height 0.25pt\vspace{6pt}
\noindent\textbf{Q.50} \hfill {\footnotesize\textsf{[NAT]}}
\par\vspace{2pt}
{\bfseries The instruction pipeline of a RISC processor has the following stages: Instruction Fetch (IF),
Instruction Decode (ID), Operand Fetch (OF), Perform Operation (PO) and Writeback (WB).
The IF, ID, OF and WB stages take 1 clock cycle each for every instruction. Consider a
sequence of 100 instructions. In the PO stage, 40 instructions take 3 clock cycles each, 35
instructions take 2 clock cycles each, and the remaining 25 instructions take 1 clock cycle
each. Assume that there are no data hazards and no control hazards.
\par\smallskip
The number of clock cycles required for completion of execution of the sequence of
instructions is \_\_\_\_\_\_. \par}
\par\vspace{3pt}\noindent\textbf{Answer:} \rule{3.2cm}{0.4pt}
\par\vspace{6pt}\hrule height 0.25pt\vspace{6pt}
\noindent\textbf{Q.51} \hfill {\footnotesize\textsf{[NAT]}}
\par\vspace{2pt}
{\bfseries A processor has 16 integer registers (R0, R1, .. , R15) and 64 floating point registers (F0,
F1,… , F63). It uses a 2-byte instruction format. There are four categories of instructions:
Type-1, Type-2, Type-3, and Type-4. Type-1 category consists of four instructions, each
with 3 integer register operands (3Rs). Type-2 category consists of eight instructions, each
with 2 floating point register operands (2Fs). Type-3 category consists of fourteen
instructions, each with one integer register operand and one floating point register operand
(1R+1F). Type-4 category consists of N instructions, each with a floating point register
operand (1F).
\par\smallskip
The maximum value of N is \_\_\_\_\_\_\_\_\_\_. \par}
\par\vspace{3pt}\noindent\textbf{Answer:} \rule{3.2cm}{0.4pt}
\par\vspace{6pt}\hrule height 0.25pt\vspace{6pt}
\noindent\textbf{Q.52} \hfill {\footnotesize\textsf{[NAT]}}
\par\vspace{2pt}
{\bfseries Given a language 𝐿, define 𝐿$i$ as follows:
𝐿0 = \{𝜀\}
𝐿$i$= 𝐿$i$−1 ⋅𝐿 𝑓𝑜$r$ 𝑎𝑙𝑙 $i$> 0
\par\smallskip
The order of a language L is defined as the smallest k such that 𝐿$k$= 𝐿$k$+1.
Consider the language L1 (over alphabet 0) accepted by the following automaton.
\par\smallskip
The order of L1 is \_\_\_\_\_.
0
0
0 \par}
\par\vspace{3pt}\noindent\textbf{Answer:} \rule{3.2cm}{0.4pt}
\par\vspace{6pt}\hrule height 0.25pt\vspace{6pt}
\noindent\textbf{Q.53} \hfill {\footnotesize\textsf{[NAT]}}
\par\vspace{2pt}
{\bfseries Consider a storage disk with 4 platters (numbered as 0, 1, 2 and 3), 200 cylinders (numbered
as 0, 1, … , 199), and 256 sectors per track (numbered as 0, 1, … , 255). The following 6
disk requests of the form [sector number, cylinder number, platter number] are received by
the disk controller at the same time:
[120, 72, 2] , [180, 134, 1] , [60, 20, 0] , [212, 86, 3] , [56, 116, 2] , [118, 16, 1]
\par\smallskip
Currently the head is positioned at sector number 100 of cylinder 80, and is moving towards
higher cylinder numbers. The average power dissipation in moving the head over 100
cylinders is 20 milliwatts and for reversing the direction of the head movement once is 15
milliwatts. Power dissipation associated with rotational latency and switching of head
between different platters is negligible.
\par\smallskip
The total power consumption in milliwatts to satisfy all of the above disk requests using the
Shortest Seek Time First disk scheduling algorithm is \_\_\_\_\_\_\_. \par}
\par\vspace{3pt}\noindent\textbf{Answer:} \rule{3.2cm}{0.4pt}
\par\vspace{6pt}\hrule height 0.25pt\vspace{6pt}
\noindent\textbf{Q.54} \hfill {\footnotesize\textsf{[NAT]}}
\par\vspace{2pt}
{\bfseries Consider an IP packet with a length of 4,500 bytes that includes a 20-byte IPv4 header and
a 40-byte TCP header. The packet is forwarded to an IPv4 router that supports a Maximum
Transmission Unit (MTU) of 600 bytes. Assume that the length of the IP header in all the
outgoing fragments of this packet is 20 bytes. Assume that the fragmentation offset value
stored in the first fragment is 0.
\par\smallskip
The fragmentation offset value stored in the third fragment is \_\_\_\_\_\_\_. \par}
\par\vspace{3pt}\noindent\textbf{Answer:} \rule{3.2cm}{0.4pt}
\par\vspace{6pt}\hrule height 0.25pt\vspace{6pt}
\end{multicols}
\vspace{8pt}
\begin{center}
\rule{0.35\linewidth}{0.5pt}\\[2pt]
\textbf{\small --- END OF QUESTION PAPER ---}\\[-8pt]
\rule{0.35\linewidth}{0.5pt}
\end{center}
\vspace{10pt}
\subsection*{OFFICIAL ANSWER KEY --- GATE 2018 (CS)}
\vspace{-4pt}
\noindent{\footnotesize\textit{Official final answer keys published by the GATE Organising Institute. For NAT questions, allowable numerical ranges are displayed where applicable. MTA = Marks To All.}}
\vspace{4pt}
\begin{center}
\small
\resizebox{\linewidth}{!}{%
\begin{tabular}{|c|c|l||c|c|l||c|c|l||c|c|l|}
\hline
\textbf{Q} & \textbf{Type} & \textbf{Key} & \textbf{Q} & \textbf{Type} & \textbf{Key} & \textbf{Q} & \textbf{Type} & \textbf{Key} & \textbf{Q} & \textbf{Type} & \textbf{Key} \\
\hline\hline
1 & MCQ & B & 9 & MCQ & A & 24 & NAT & 2 & 40 & MCQ & C \\ \hline
1 & MCQ & D & 10 & MCQ & B & 25 & NAT & 34 & 41 & MCQ & C \\ \hline
2 & MCQ & A & 10 & MCQ & C & 26 & MCQ & D & 42 & MCQ & B \\ \hline
2 & MCQ & A & 11 & MCQ & A & 27 & MCQ & D & 43 & NAT & 109 \\ \hline
3 & MCQ & D & 12 & MCQ & D & 28 & MCQ & A & 44 & NAT & 0.60 to 0.62 \\ \hline
4 & MCQ & C & 13 & MCQ & C & 29 & MCQ & A & 45 & NAT & 10230 \\ \hline
4 & MCQ & D & 14 & MCQ & C & 30 & MCQ & A & 46 & NAT & 80 \\ \hline
5 & MCQ & B & 15 & NAT & 0.021 to 0.024 & 31 & MCQ & C & 48 & NAT & 16 \\ \hline
5 & MCQ & D & 16 & NAT & 0.27 to 0.30 & 32 & MCQ & B & 50 & NAT & 219 \\ \hline
6 & MCQ & D & 17 & NAT & 3 & 33 & MCQ & C & 51 & NAT & 32 \\ \hline
6 & MCQ & B & 18 & NAT & 3 & 34 & MCQ & B & 52 & NAT & 2 \\ \hline
7 & MCQ & B & 19 & NAT & 42 & 35 & MCQ & B & 53 & NAT & 85 \\ \hline
7 & MCQ & C & 20 & NAT & 4 & 36 & MCQ & D & 54 & NAT & 144 \\ \hline
8 & MCQ & B & 21 & NAT & 4 & 37 & MCQ & D &  & &  \\ \hline
8 & MCQ & B & 23 & NAT & 59.0 to 60.0 & 39 & MCQ & A &  & &  \\ \hline
\end{tabular}%
}
\end{center}
\end{document}